\section{CROSS-MODALITY EVIDENCE SYNTHESIS}
\label{sec:full_corpus_synthesis}

The primary synthesis comprised 8,203 governed claims: 3,020 evidence
statements, 4,779 metric results, and 404 trade-off records. The broader set of
8,234 non-quarantined claims additionally contained 31 context-only metric
records; these records were excluded from primary prevalence and performance
summaries. A conditional subset of 1,505 primary claims met a recorded
comparability or guardrail condition and is reported as a subset rather than
as unconditional evidence.

\subsection{S1: Optical Modality Taxonomy}

Photonic-terahertz systems formed the largest mutually exclusive modality
group (69/206, 33.5\%), followed by fiber systems (56, 27.2\%), VLC/LiFi (38,
18.4\%), FSO (31, 15.0\%), hybrid optical systems (9, 4.4\%), and other optical
platforms (3, 1.5\%). The distribution shows that O-ISAC is not a single-medium
literature: photonic-terahertz and fiber together accounted for 60.7\% of the
included studies, while optical wireless modalities remained a substantial
and technically distinct evidence base.

\begin{table}[!t]
\caption{Mutually exclusive optical-modality distribution.}
\label{tab:s1_modality}
\centering
\begin{tabular}{lrr}
\toprule
\textbf{Modality} & \textbf{Studies} & \textbf{\% of 206} \\
\midrule
Photonic-THz & 69 & 33.5 \\
Fiber & 56 & 27.2 \\
VLC/LiFi & 38 & 18.4 \\
FSO & 31 & 15.0 \\
Hybrid optical & 9 & 4.4 \\
Other optical & 3 & 1.5 \\
\bottomrule
\end{tabular}
\end{table}

\subsection{S2: Architecture and Integration Mechanisms}

Integration was predominantly resource- and hardware-centered. Shared
resource allocation appeared in 118 studies (57.3\%), shared hardware in 117
(56.8\%), and a shared waveform in 113 (54.9\%). Shared links or channels were
reported in 87 studies (42.2\%), joint design or optimization in 72 (35.0\%),
shared optical carriers in 49 (23.8\%), and a shared application scenario in
46 (22.3\%). These categories were multi-label and therefore represent
overlapping design mechanisms rather than a partition of the corpus.

The pattern indicates that the dominant O-ISAC integration logic is not merely
co-location of two functions. More than half of the studies reused hardware,
waveform structure, or resource allocation, while over one-third explicitly
introduced joint design or optimization. Shared-application-only integration
was less common and was kept separate from stronger physical or algorithmic
coupling.

\begin{table*}[!t]
\caption{Multi-label architecture and integration mechanisms.}
\label{tab:s2_integration}
\centering
\begin{tabular}{lrr@{\hspace{2em}}lrr}
\toprule
\textbf{Mechanism} & \textbf{$n$} & \textbf{\%} &
\textbf{Mechanism} & \textbf{$n$} & \textbf{\%} \\
\midrule
Shared resource allocation & 118 & 57.3 & Shared hardware & 117 & 56.8 \\
Shared waveform & 113 & 54.9 & Shared link/channel & 87 & 42.2 \\
Joint design/optimization & 72 & 35.0 & Shared optical carrier & 49 & 23.8 \\
Shared application scenario & 46 & 22.3 & Mixed mechanism & 3 & 1.5 \\
\bottomrule
\end{tabular}
\end{table*}

\subsection{S3: Communication and Sensing Metric Reporting}

The metric map contained 4,779 primary metric claims after excluding 31
context-only and 51 quarantined metric claims. Sensing was the largest coded
domain, with 1,816 claims across 203 studies. Communication accounted for
1,328 claims across 194 studies, joint metrics for 870 claims across 158
studies, and implementation metrics for 476 claims across 64 studies. The
remaining 289 primary claims were distributed across smaller system, energy,
scenario, waveform, channel, security, machine-learning, component, network,
and other source-defined domains.

\begin{table}[!t]
\caption{Largest primary metric domains.}
\label{tab:s3_metric_domains}
\centering
\begin{tabular}{lrrr}
\toprule
\textbf{Domain} & \textbf{Studies} & \textbf{Claims} & \textbf{Conditional} \\
\midrule
Sensing & 203 & 1,816 & 505 \\
Communication & 194 & 1,328 & 296 \\
Joint & 158 & 870 & 260 \\
Implementation & 64 & 476 & 28 \\
All other domains & --- & 289 & --- \\
\midrule
Total & --- & 4,779 & --- \\
\bottomrule
\end{tabular}
\end{table}

The breadth of the source-specific metric vocabulary precluded a defensible
single performance ranking. Similar labels frequently referred to different
roles, planes, scenarios, or validation contexts. Accordingly, the metric map
supports reporting-pattern and within-stratum comparisons, but not a universal
O-ISAC performance league table.

\subsection{S4: Communication--Sensing Trade-offs}

Trade-off evidence comprised 404 primary records contributed by 168 studies.
The most frequently reported family concerned bandwidth, spectrum, or resource
allocation (95 claims; 71 studies), followed by power, energy, or dynamic range
(90; 64), communication reliability versus sensing quality (59; 48),
rate--resolution (46; 39), and rate--accuracy or localization (40; 33). Less
frequent families concerned waveform/hardware complexity (21; 20),
rate--range/coverage (21; 16), other joint trade-offs (19; 15), and qualitative
or partial relationships (11; 9). One security/resilience trade-off and one
non-antagonistic synergy relationship were also retained.

\begin{table*}[!t]
\caption{Communication--sensing trade-off families.}
\label{tab:s4_tradeoffs}
\centering
\begin{tabular}{lrrr}
\toprule
\textbf{Trade-off family} & \textbf{Studies} & \textbf{Claims} & \textbf{Conditional subset} \\
\midrule
Bandwidth/spectrum/resource allocation & 71 & 95 & 90 \\
Power/energy/dynamic range & 64 & 90 & 81 \\
Communication reliability vs. sensing quality & 48 & 59 & 55 \\
Rate--resolution & 39 & 46 & 41 \\
Rate--accuracy/localization & 33 & 40 & 39 \\
Waveform/hardware/complexity & 20 & 21 & 21 \\
Rate--range/coverage & 16 & 21 & 18 \\
Other joint trade-off & 15 & 19 & 18 \\
Qualitative/partial general relation & 9 & 11 & 8 \\
Security/resilience & 1 & 1 & 1 \\
Synergy/non-antagonistic coupling & 1 & 1 & 1 \\
\bottomrule
\end{tabular}
\end{table*}

These results broaden the conventional ``rate versus sensing'' narrative.
Resource allocation and power constraints were at least as prominent as
rate-centered formulations, suggesting that practical O-ISAC design should be
described as a multi-resource optimization problem rather than by one Pareto
axis alone.

\subsection{S5: Validation Maturity and Benchmark Readiness}

Laboratory experiment or proof-of-concept evidence was the most common maximum
validation tier (78 studies, 37.9\%), followed by controlled prototypes (66,
32.0\%). Enhanced simulation or dataset-supported validation accounted for 18
studies (8.7\%), basic simulation or numerical evidence for 32 (15.5\%), and
field or deployment-oriented validation for only 12 (5.8\%). Thus, 156 studies
(75.7\%) reached at least a laboratory/proof-of-concept tier, but only a small
minority extended to field conditions.

Artifact availability was substantially weaker. Data were openly available
for 13 studies (6.3\%) and available on request for 41 (19.9\%); 145 (70.4\%)
reported no available dataset or did not report availability, and 7 were not
applicable. Code or model artifacts were unavailable or not reported in 197
studies (95.6\%); 7 offered access on request, 1 exposed partial components,
and 1 was not applicable. This gap between experimental maturity and artifact
availability explains why benchmark readiness remained limited despite strong
within-study reporting in much of the corpus.

\subsection{S6: Enabling Technologies and Applications}

The enabling-technology map was multi-label. Photonic-terahertz or
microwave-photonic generation appeared in 68 studies (33.0\%), FMCW/chirped
techniques in 66 (32.0\%), coherent optical techniques in 64 (31.1\%), and
OFDM or closely related multicarrier methods in 56 (27.2\%). Fiber distributed
acoustic sensing appeared in 22 studies (10.7\%); photonic integration and
ML/AI appeared in 20 studies each (9.7\%). Beamforming appeared in 13 studies,
optical phased arrays in 11, MIMO in 7, and RIS/ORIS and digital-twin methods
in 2 each. Nineteen studies (9.2\%) had no recognized broad technology label
and therefore remained in the fallback category.

Application coding likewise allowed overlap. Explicit 6G/access-network use
appeared in 100 studies (48.5\%), environmental monitoring in 55 (26.7\%),
vehicular applications and smart infrastructure in 41 each (19.9\%),
industrial settings in 34 (16.5\%), optical access networks in 32 (15.5\%),
security/surveillance in 26 (12.6\%), and indoor positioning in 25 (12.1\%).
Aerospace applications appeared in 20 studies, healthcare in 8, datacenter
contexts in 4, and underwater settings in 2. Fifteen studies (7.3\%) had no
recognized broad application label. Long-tail or unmatched source terms were
retained in the normalization audit; a fallback category was used only when a
study had no recognized category on the relevant axis.

\subsection{S7: Evidence Gaps and 6G Relevance}

The frozen study-level crosswalk classified 138 studies (67.0\%) as directly
6G-relevant and 64 (31.1\%) as inferentially relevant. Only one study was weakly
relevant and three were not applicable. This distribution supports O-ISAC as a
substantive 6G research direction, while the appraisal and availability results
show that thematic relevance has advanced faster than benchmarking maturity.

The principal research needs are therefore: (i) explicit measurement-plane and
metric-role contracts; (ii) scenario-complete reporting of operating
conditions; (iii) shared, externally usable baselines; (iv) open data, code,
models, and experimental protocols; (v) field validation beyond controlled
laboratory demonstrations; and (vi) modality-aware benchmarks that do not
collapse physically different sensing constructs into nominally similar metric
labels. These priorities are supported jointly by the metric-governance layer,
the trade-off map, the TQAF appraisal, and the artifact-availability results.

\subsection{Interpretive Boundary}

The reported distributions characterize the governed 206-study corpus. They
do not imply that every claim was independently verified by a second human
reviewer, and they do not establish pooled causal effects. The appropriate
interpretation is a transparent, investigator-supervised, AI-assisted, and
claim-governed synthesis whose most reproducible conclusions concern the
structure, reporting practices, validation maturity, and open benchmarking
needs of the field.
